\section{引言}

量子色动力学（QCD）把强子物质与其基本组分——夸克和胶子——联系起来。尽管经历了四十余年紧张的
实验与理论研究，QCD 的许多方面仍不甚明了。作为这项努力的一部分，轻子对核子的深度非弹性散射（DIS）
在确立 QCD 作为强相互作用主宰理论的过程中发挥了核心作用，并且至今仍是为解开 QCD 之谜所做的种种
尝试的重要支柱（完整的综述可参见，例如，Refs.~\cite{pdg12} 与
\cite{Collins:2011zzd}）。

在理论上，深度非弹性散射的包含反应截面被分离为轻子张量与强子张量，二者分别刻画散射过程中的电弱
动力学与强动力学。强子张量因子化为红外安全的微扰系数与部分子分布函数（PDF）的卷积；PDF 承载
低能 QCD 物理，因此必须以非微扰方式确定。PDF 与散射过程无关——但依赖于靶核子，而微扰系数则与
外部散射态无关。

PDF 之所以重要有两个原因。其一，它以强相互作用基本理论的语言直接给出强子态——宇宙中可见物质的
主要形态——之组分的知识。其二，它为对撞机实验中的强子本底提供约束，例如大型强子对撞机上的实验。
这些本底影响多种高能实验的灵敏度，包括希格斯玻色子性质的研究
\cite{Thorne:2011kq,Alekhin:2010dd,Alekhin:2011ey}，以及重 $W'$
和 $Z'$ 玻色子产生的寻找 \cite{Brady:2011hb}。

目前尚不能借助第一性原理的格点 QCD 直接从 QCD 计算 PDF，因为 PDF 由光锥矩阵元定义，
而这些矩阵元无法在欧氏格点上直接获取。因此，PDF 通常由对不同实验的整体分析提取
\cite{Jimenez-Delgado:2014xza,Jimenez-Delgado:2013boa,Ball:2013lla,
Arbabifar:2013tma,Owens:2012bv,Accardi:2011fa,Leader:2010rb,Blumlein:2010rn,
Hirai:2008aj}。

自然地，人们希望对 PDF 进行直接的非微扰计算：只要精度足够，这样的计算将进一步约束整体拟合中
实验上可能难以到达的区域。

直到最近，格点计算主要集中于确定 PDF 的 Mellin 矩；这些矩与 twist-2 算符的矩阵元相关，
而后者\emph{可以}在格点上确定
\cite{Bali:2012av,Braun:2008ur,Guagnelli:2004ga}。格点调节子破坏洛伦兹对称性，
从而在不同质量维算符之间诱导辐射混合。由此产生的格点间距的幂次发散完全掩盖了连续极限。
通过仔细选取外部动量与算符，可以提取到四阶矩，但这显著限制了从 PDF 提取有意义结果的精度
\cite{Detmold:2003rq,Detmold:2001dv}。
超出四阶矩后，幂次发散混合不可避免，该方法便无法提供可靠的计算结果。

Ji 近来提出一种从格点计算中直接提取 PDF 的新途径 \cite{Ma:2014jla,Ji:2013dva}，
其基础是大动量有效理论 \cite{Ji:2014gla}。初步结果最早见于 Ref.~\cite{Lin:2014zya}，
但仍存在若干问题，包括对相关格点矩阵元重正化的完整理解，以及在格点上实际分辨足够大动量的能力。

本文提出一种新框架，它在连续极限下消除混合，并且在原则上使从格点 QCD 提取 PDF 更高阶矩成为可能。
我们将此框架称为``涂抹算符乘积展开''（smeared operator product expansion，sOPE）。我们把连续统
矩阵元按局域涂抹的格点自由度基底展开。所得矩阵元是两个尺度的函数：涂抹尺度与重正化尺度。sOPE 提供
了必要的框架，可借助威尔逊系数把这些矩阵元与 DIS 强子张量等在唯象上有用的量联系起来。

涂抹技术在格点 QCD 中被广泛用于压低紫外涨落、部分恢复转动对称性，从而系统地提高格点计算的精度
\cite{Morningstar:2003gk,Hasenfratz:2001hp,Bernard:1999kc,Albanese:1987as}。关于格点计算中
涂抹的教学性概述可参见，例如，Ref.~\cite{Gattringer:2010zz}。在 sOPE 中，我们经由梯度流实现涂抹；
梯度流是场在一个新维度——流时间——中朝作用量驻点的经典演化
\cite{Luscher:2013cpa,Luscher:2011bx,Lohmayer:2011si,Luscher:2010iy,
Narayanan:2006rf}。在格点上非微扰确定的矩阵元无须进一步重正化（费米子波函数重正化除外
\cite{Luscher:2013cpa}），并且只要在连续极限下涂抹尺度以物理单位保持固定便保持有限
\cite{Makino:2014sta,Luscher:2011bx}。这还有两个优点：梯度流允许只使用一两个格距的涂抹长度，
远小于典型格点上的强子长度尺度，因而不会扭曲低能物理 \cite{Davoudi:2012ya}；此外梯度流的计算代价
非常低廉。

梯度流如今已作为研究格点规范理论的成熟工具确立下来，应用范围涵盖从标度设定（即以物理单位确定格距）
\cite{Borsanyi:2012zr,Borsanyi:2012zs,Bazavov:2013gca,Fodor:2014cpa,
Cheng:2014jba}，到研究格点规范理论中的重正化。例如，梯度流已被用于为强耦合常数定义有限体积重正化方案
\cite{Fodor:2012td,Fritzsch:2013yxa,Fritzsch:2013hda,
Ramos:2013gda,Rantaharju:2013bva}，用于算符重正化 \cite{Monahan:2013lwa}，
以及理解重正化矩阵元的非微扰标度依赖 \cite{Luscher:2014kea}。
相关工作还利用 Ward 恒等式研究了格点上的手性费米子 \cite{Shindler:2013bia}
与能动量张量 \cite{Suzuki:2013gza,DelDebbio:2013zaa, Makino:2014taa}。

本文对具有四次型相互作用的单个实标量场引入梯度流，并概述应用于实标量场论的 sOPE 的若干性质。
标量场论描述自然界中某些重要的临界现象，例如反铁磁相变，并且作为四维量子场论基本思想的试金石有着
悠久的历史。就本文的目的而言，欧氏标量场论的主要优点在于其简单性：它使 sOPE 的结构一目了然，
并凸显其与局域算符乘积展开（OPE）的相似与差异之处。

下一节我们回顾 Wilson 的 OPE 并引入 sOPE。随后在第~\ref{sec:phi4}~节把 sOPE 应用于标量场论，
在第~\ref{sec:lattsmear}~节说明梯度流如何消除幂次发散混合，并在第~\ref{sec:wc}~节把微扰威尔逊系数
计算到两圈。最后在第~\ref{sec:rgeqns}~节，我们通过威尔逊系数的重正化群方程研究 sOPE 的标度依赖。
